\documentclass{book}
\usepackage{amsmath,amssymb,amsthm}
\usepackage{hyperref}

\newtheorem{theorem}{Theorem}[section]
\newtheorem{proposition}[theorem]{Proposition}
\newtheorem{corollary}[theorem]{Corollary}
\theoremstyle{definition}
\newtheorem{definition}{Definition}[section]
\newtheorem{example}{Example}[section]
\theoremstyle{remark}
\newtheorem*{remark}{Remark}

\title{Lecture Notes on Real Analysis}
\author{Prof. Frank Miller}

\begin{document}
\maketitle
\tableofcontents

\chapter{Sequences and Limits}

\section{Convergence}

\begin{definition}
A sequence $(a_n)$ of real numbers converges to $L \in \mathbb{R}$ if for every $\varepsilon > 0$ there exists $N \in \mathbb{N}$ such that $|a_n - L| < \varepsilon$ whenever $n \geq N$.
\end{definition}

\begin{theorem}[Uniqueness of limits]
If $a_n \to L$ and $a_n \to M$, then $L = M$.
\end{theorem}

\begin{proof}
Suppose $L \neq M$. Take $\varepsilon = |L - M|/3$. Both definitions give $N$ past which $|a_n - L| < \varepsilon$ and $|a_n - M| < \varepsilon$. Then $|L - M| \leq |L - a_n| + |a_n - M| < 2\varepsilon$, contradiction.
\end{proof}

\begin{example}
The sequence $a_n = 1/n$ converges to $0$. Given $\varepsilon$, pick $N > 1/\varepsilon$.
\end{example}

\begin{remark}
Convergence is a local property: changing finitely many terms cannot affect the limit.
\end{remark}

\section{Cauchy Sequences}

\begin{definition}
A sequence $(a_n)$ is Cauchy if for every $\varepsilon > 0$ there exists $N$ such that $|a_m - a_n| < \varepsilon$ for all $m, n \geq N$.
\end{definition}

\begin{proposition}
Every convergent sequence is Cauchy.
\end{proposition}
\end{document}
